\documentclass[10pt,a4paper]{article}
\usepackage[latin1]{inputenc}
\usepackage{amsmath}
\usepackage{amsfonts}
\usepackage{amssymb}
\usepackage{graphicx}
\usepackage{subcaption}
\usepackage{lmodern}
\usepackage{placeins}
\usepackage{multicol}
\usepackage{tikz}
\usetikzlibrary{shapes.arrows,chains, positioning}
\usepackage{booktabs} %professional tables
\usepackage[left=2cm,right=2cm,top=2cm,bottom=2cm]{geometry}
\usepackage{microtype} %optimises spacing, needs to go after fonts
\usepackage{hyperref} %adds links (also in TOC), should be loaded at the very end
\usepackage{caption}
\usepackage{longtable}

\author{Thomas Keck\\
Christian Pulvermacher}
\begin{document}

%\setlength{\textfloatsep}{5pt plus 1.0pt minus 1.0pt}
%\setlength{\floatsep}{5pt plus 1.0pt minus 1.0pt}
%\setlength{\intextsep}{5pt plus 1.0pt minus 1.0pt}

\date{\today}
\title{Full Event Interpretation Report}

\maketitle

\begin{abstract}
This report contains key performance indicators and control plots of the Full Event Interpretation algorithm,
for a training on ${NEvents} events.
The user-, pre-, and post-cuts as well as trained multivariate selection methods are described.
Furthermore the resulting purities and efficiencies are stated. In ${overallSignalEfficiencyInPercent}\,\% of the events
the Full Event Interpretation reconstructed a final particle (${finalParticleNames}) correctly.

\begin{center}
\includegraphics[width=0.5\textwidth]{nordbert.pdf}

\end{center}
\end{abstract}


\clearpage

\tableofcontents

\FloatBarrier
\clearpage
\section{Summary}
For each final-state particle a multivariate selection method is trained without any previous cuts
on the candidates. Afterwards, a post cut is applied on the signal probability calculated by the method.
This reduces combinatorics in the following stages of the Full
Event Interpretation.
\begin{center}
\begin{table}[h]
\caption{Final-state particle efficiency and purity before and after the applied post-cut.}
\centering
\begin{tabular}{c|rr|rr}
\toprule
Final-state &  \multicolumn{2}{c}{Efficiency in \%}  &  \multicolumn{2}{c}{Purity in \%} \\
particle    &  detector & post-cut   &  detector & post-cut \\
\midrule
${finalStateParticleEPTable}
\bottomrule
\end{tabular}
\end{table}
\end{center}
\FloatBarrier
For each decay channel of each intermediate particle a multivariate selection method is trained after applying a fast pre-cut
on the candidates. Afterwards, a post-cut is applied on the signal probability calculated by the method.
This reduces combinatorics in the following stages of the Full
Event Interpretation. For some intermediate particles, in particular the final B mesons, an additional user-cut
is applied before all other cuts.
\begin{center}
\begin{table}[h]
\caption{Per-particle efficiency and purity before and after the applied user-, pre- and post-cut.}
\centering
\begin{tabular}{c|rrrr|rrrr}
\toprule
Particle &  \multicolumn{4}{c}{Efficiency in \%}  &  \multicolumn{4}{c}{Purity in \%} \\
         &  reconstruction & user-cut & pre-cut & post-cut   & reconstruction & user-cut & pre-cut & post-cut \\
\midrule
${combinedParticleEPTable}
\bottomrule
\end{tabular}
\end{table}
\end{center}

${finalParticleInputs}

\clearpage
${cpuTimeSummary}

\FloatBarrier
\clearpage
\section{Particle configuration}
%\begin{multicols}{2}
\begin{verbatim}
${particleConfigurations}
\end{verbatim}
%\end{multicols}
\FloatBarrier
\clearpage
\section{Final state particles}
${finalStateParticleInputs}

\FloatBarrier
\clearpage
${combinedParticleInputs}

\end{document}
